% Build twice with XeLaTeX.
\documentclass[11pt,a4paper]{article}
\usepackage[margin=21mm,headheight=15pt]{geometry}
\usepackage{fontspec,xcolor,tabularx,booktabs,fancyhdr,hyperref}
\IfFontExistsTF{TeX Gyre Pagella}{\setmainfont{TeX Gyre Pagella}}{\setmainfont{Liberation Serif}}
\IfFontExistsTF{TeX Gyre Heros}{\setsansfont{TeX Gyre Heros}}{\setsansfont{Liberation Sans}}
\setmonofont{Liberation Mono}
\definecolor{ink}{HTML}{182021}
\definecolor{teal}{HTML}{0E6059}
\hypersetup{colorlinks=true,linkcolor=teal,urlcolor=teal,pdftitle={EDAgent Project Plan v4}}
\pagestyle{fancy}
\fancyhf{}
\fancyhead[L]{\small\sffamily EDAgent \enspace / \enspace Project Plan v4}
\fancyhead[R]{\small\sffamily 14 November 2026}
\fancyfoot[L]{\small\sffamily Allegheny County road-safety investigation}
\fancyfoot[R]{\small\sffamily\thepage}
\renewcommand{\headrulewidth}{.35pt}
\setlength{\parindent}{0pt}
\setlength{\parskip}{5pt}
\setlength{\emergencystretch}{2em}
\renewcommand{\arraystretch}{1.22}
\newcolumntype{Y}{>{\raggedright\arraybackslash}X}
\newcommand{\code}[1]{{\Urlmuskip=0mu plus 1mu\relax\nolinkurl{#1}}}
\newcommand{\pagehead}[2]{\clearpage\phantomsection\pdfbookmark[0]{#2}{page.\thepage}{\sffamily\small\color{teal}\MakeUppercase{#1}}\par{\LARGE\raggedright\bfseries #2}\par\vspace{3pt}}
\newcommand{\sub}[1]{\par\smallskip{\large\bfseries #1}\par}
\newcommand{\lead}[2]{\textbf{#1}\enspace #2\par}
\newcommand{\lib}[2]{\href{https://chatgpt.com/library?file_id=#1\&account_id=personal}{#2}}
\newcommand{\repo}[2]{\href{https://github.com/ExploratoryDataAnalysis/EDAgent/#1}{#2}}
\begin{document}
\color{ink}
{\sffamily\small\color{teal} PROJECT CONTINUATION \quad / \quad 14 NOVEMBER 2026}\par
{\Huge\bfseries EDAgent Project Plan v4}\par
{\Large Public investigation and release verification}\par
\vspace{4pt}\hrule\vspace{5pt}
\sub{Phase 6 handoff}
Phase 6 delivered the public \code{/road-safety} shell and verified aggregate API, chart and table alternatives, the separate \code{/illustrative-scenario} workflow, and Compose/PostgreSQL staging. The signed-in general analyser retained its own dataset \code{/overview}, upload, analysis and export. These capabilities were exercised on the Phase 6 build. The current deployment and later \code{main} require fresh verification.

{\small
\begin{tabularx}{\linewidth}{@{}p{29mm}Y@{}}
\toprule Evidence date & Disposition and boundary\\\midrule
13 Nov 2026 & Commit \repo{commit/bcdf44d}{bcdf44d}, from \code{71376b3}, includes D7 contract lineage \code{6d01d2a} and Phase 5 arithmetic \code{0d136f3}. PR 11 merged as \repo{commit/87a7aea}{87a7aea}, the audit baseline.\\
13 Nov 2026 & Handoff records 9 real-candidate artifact checks, 140 road-safety checks with 9 historical skips, 82 focused API/intervention checks and 192 backend checks passed. Frontend build passed with a chunk-size warning. Chromium synthetic integration passed at 1440 and 390 px, including keyboard, chart alternatives, null/zero, stale responses and real-site refusal.\\
13 Nov 2026 & Loopback HTTP staging confirmed Compose, PostgreSQL 16 at Alembic \code{0001} head with no pending schema operations, mounted hash, calculator execution, signed-in CSV upload through PDF download, and real Overview HTTP 503 with null data. Staging used the mock provider.\\
13--14 Nov 2026 & Localhost captures show shared dark workspace chrome, an unavailable Overview with retry, and an empty calculator with reference and sensitivity fields. “Sign in to analyser” is the general analyser entry. Captures establish visible states only.\\
14 Nov 2026 & The audit recorded pending external approval. GitHub \code{main} resolved to \repo{commit/93991e5}{93991e5}, PR 12. Files and diffs were inspected by revision. Current local runtime, external ledger, intended deployment and provider remain unverified. See page \pageref{baseline}.\\
\bottomrule
\end{tabularx}}
\sub{Release position}
The real \code{v3-d7-2} candidate (schema 1) reproduced to the same manifest hash in both handoff builds. Its scientific disposition remains pending on supplied evidence. A synthetic ledger exercises the gate on invented artifacts. Genuine approval requires a named, dated decision on the exact real manifest. External HTTPS and final sign-off remain open.

\textbf{Document scope} \enspace Plan only. The work packages below are proposed. The canonical application root is \code{EDAgent_Completed_Extension/EDAgent/}. Unqualified frontend, backend and Compose paths are relative to it. \lib{libfile_09c025eb71308191a73c321da0f26318}{Plan v3 source} · \lib{libfile_919f39c5968c8191b2b6ac25646f7eae4}{Plan v3 PDF} · \repo{blob/bcdf44d/docs/phase6_handoff.md}{Handoff and test commands}.

\pagehead{Governing decisions}{Scope and exclusions}
{\small
\begin{tabularx}{\linewidth}{@{}p{32mm}Y@{}}
\toprule Area & Rule for v4\\\midrule
Frozen record & WPRDC cumulative Allegheny County reported crashes, 2004--2024, resource \code{2c13021f-74a9-4289-a1e5-fe0472c89881}. Exclude the separate 2025 data. Preserve the executed v6 notebook, source and run manifests, single canonical notebook and original exports. Work from held inputs and document unavailable evidence.\\
Public evidence & Reviewed historical county/year/era and clipped grid aggregates at approved grains. Count reports in the extract. Travel exposure, per-trip risk and completeness beyond the extract are outside the claim. County GEOID 42003. Separate 2004--2015 and 2016--2024 severity semantics. Unsupported pooled serious/fatal values stay null. Preserve missing, unknown, unavailable and observed zero.\\
Attribution and exposure & Zero accepted primary crash-to-site assignments. Police-text coordinates and some coding harmonisations remain unverified. Station 375 is internal. Its observed 2023 crash count is unknown. The 11,072 vehicles/day value describes the archived 2023 traffic section containing segment 0080. No site ranking, historical site screening, site-year frequency, calibrated SPF, Empirical Bayes, short-horizon forecast or safest-route recommendation.\\
Treatments & Six references retain \code{REFERENCE_ONLY}. Calculator uses hypothetical counts and compatible factors under schema 1. Real-site input returns \code{REAL_SITE_BENEFIT_UNAVAILABLE} with no numerical benefit. Actual causal recommendations and local efficacy claims are excluded.\\
Models & Retrospective severity conditional on reported crashes and known outcomes. Preserve \code{NOT_APPROVED_FOR_PRODUCTION}, post-event feature timing and previously consulted holdout. Live scoring, public estimators, per-crash probabilities and unsupported comparisons are excluded.\\
Workspace & General-purpose analyser for user or sample data, with its own scope and ownership. D7 stays within its release contract. Provider features are optional and stay in the analyser.\\
\bottomrule
\end{tabularx}}
\sub{Science decisions outside routine implementation}
New years, changed inclusion or severity mappings, new geographies or denominators, point-level publication, new adjusted or joint distributions, and model derivatives require review of the changed evidence and exact artifact. Reopening any site, exposure, forecast or treatment no-go requires a new science decision with qualifying evidence. Record owner, dated decision, source hashes, affected claims, artifact, gate and withdrawal procedure.
\sub{What changes from v3}
V3's release certification becomes a frontend rebuild with a useful gate-closed release and a gated annual slice. V4 defines reader routes and an evidence-to-interface contract, then gives map, calculator, analyser, provider and retrospective review independent acceptance decisions. Carry v3 D0--D5 and D7--D11 forward as evidence and release checks. D6 stays \code{OUT_OF_SCOPE_NO_GO}. D12--D15 stay optional. Full research and diagnostic lineage is in v3.

\pagehead{Sequence and dependencies}{Milestones and release decisions}
{\small
\begin{tabularx}{\linewidth}{@{}p{13mm}p{15mm}p{45mm}Y@{}}
\toprule Phase & A--H & Start condition & Independent exit / owner\\\midrule
7A & A B C E & Begin deployment/D7 diagnosis, reader validation and fixture frontend in parallel. & Editorial root, standalone Methods and useful deep link. Frontend, QA, deployment.\\
7B & A C E & Build with fixtures during 7A. Promote with approval and intended HTTPS evidence. & Verified annual record and selected-year coverage. Data steward, reviewer, deployment.\\
8 & D E & Phase 7 contract and figure pattern. & Map and evidence pages independently clear artifact, parity and reader checks. Science, frontend.\\
9 & D E & Initial 7A reader tasks. Calculator contract exists. & Invented example can ship early with verified calculator route. Frontend, QA.\\
10 & F & Route split and verified auth boundary. Real sample depends on D7. & Onboarding/sample first. Private guest after lifecycle tests. Workspace, privacy.\\
11 & G & Signed-in diagnostics can start early. Guest AI depends on Phase 10. & Provider reliability first, BYOK after a separate decision. AI, privacy.\\
12 & H & Offline reproduction can start during Phase 7. & Offline, synthetic and real diagnostics have independent exits. ML, reviewer.\\
\bottomrule
\end{tabularx}}
\sub{Parallel paths and dependencies}
{\small
\begin{tabularx}{\linewidth}{@{}p{30mm}Y@{}}
\toprule Track & Progress and exposure criteria\\\midrule
Editorial & Reviewed scope and Methods, reader prototypes and shell PR → fixture/browser checks → 7A on a verified safe deployment. D7 approval may remain pending.\\
Data and hosting & Pin revision and runtime → re-establish artifact gate and auth boundary → genuine decision and target HTTPS rehearsal → 7B. Review and hosting verification can run concurrently.\\
Integration & While hosting is pending, finish the shell and fixture tests locally and leave external verification open. While approval is pending, retain the independently useful 7A layout.\\
Independent tracks & Calculator can clear before Phase 8 or 7B. Map and severity follow the first evidence slice. Provider, BYOK, guest upload and model lab have independent release paths.\\
\bottomrule
\end{tabularx}}
\sub{First engineering session}
Assign deployment owner, reviewer, frontend lead and QA. Pin audited, inspected and target builds. Address post-handoff gate and auth differences. Open concurrent work items for the 7A shell, fixture contract and baseline captures, and reader tasks. Review the first PR on page \pageref{firstpr} against the gate-closed route contract.
\pagehead{Phase 7 · baseline and release contract}{Phase 7 — Public investigation frontend and verified release}
\label{baseline}
\lead{Objective and owners}{Frontend and science editor make the public investigation useful on its own. Deployment maintainer and data steward diagnose the runtime and candidate in parallel. QA owns reproducible acceptance and the deployment owner signs the exposure decision. Missing evidence and weak presentation are separate problems.}
\sub{Pin the deployment and restore the release boundary}
\lead{Observed on 14 November}{Audit and implementation-plan findings describe \code{87a7aea}. Static inspection of later \code{93991e5} found that \code{reviewed_release} returns a canonical fallback when the approvals setting is absent, and the fallback presents approved-like provenance. Added \code{/spatial/hotspots} and \code{/models/comparison} under the Overview API use canonical data outside the Phase 6 dependency. The merge also adds guest/open access and an unauthenticated identity fallback that can reuse the earliest active user. These are source findings. The screenshots and handoff do not establish that this code was running on localhost or any public host.}
\lead{Release blocker}{Use the Phase 6 gates as the contract for the chosen target revision. Restore default denial, exact-byte provenance and strict protected workspace ownership. Remove or close ungated legacy/canonical aggregates, model and hotspot routes and review client fallback/download paths. Keep only intended anonymous routes. Test protected reads without an identity and cross-owner isolation with harmless fixtures. Verify safe migrations and retained account ownership. Use an independently verified earlier image or isolate affected surfaces if the target cannot be corrected in time.}
\lead{Route inventory}{At \code{93991e5}, \code{frontend/src/App.tsx} mounts all routes under \code{AppLayout}: root, \code{/road-safety}, \code{/illustrative-scenario}, \code{/login}, \code{/upload}, \code{/overview}, \code{/inventory}, \code{/workspace}, \code{/chat}, \code{/visualize}, \code{/history}, \code{/reports}, \code{/exports}; unmatched paths redirect to root. \code{/methods} and \code{/analyser} are proposed. \code{backend/app/api/v1/router.py} mounts Overview, interventions, legacy frequency, auth, datasets, profiling, analysis, chat, sessions, reports, exports, jobs and health under API v1. Validate the running route inventory independently.}
\lead{Diagnostic order}{Identify Docker host, target origin, code/image IDs, Compose layers, database and migration head. Check effective proxy and API base URL in the browser, ingress/TLS and direct SPA refresh. Check API health, container mount and config presence, provenance response/body and server-side gate reason. Look for missing/wrong ledger, candidate mismatch, changed/missing artifacts, unsupported schema, API routing or transport failure. Record a redacted request trace, decision and reproduction. The live host, ledger, exact mounted bytes and provider are access gaps: deployment owner verifies the first three here. Phase 11 owns the provider test.}
\lead{D7 identity}{Schema 1, \code{descriptive_release}, \code{v3-d7-2}. Compare the mounted manifest and both archived candidate builds to the handoff SHA-256, then verify each declared output. Current review and hosting decisions are pending.}
\begin{center}\small\code{ec9def348fc67561f7ee2a5c7b3307a7db42b6e8832d7df7cc04252b87e70a82}\end{center}
\lead{Deployment inputs}{Use \code{docker-compose.yml} and \code{docker-compose.staging.yml} from the canonical application as starting points. Map host \code{D7_RELEASE_ROOT} and \code{D7_APPROVAL_LEDGER} to read-only in-container \code{ROAD_SAFETY_RELEASE_ROOT} and \code{ROAD_SAFETY_APPROVALS}; verify effective configuration. Keep the ledger external and isolate credentials and storage. Confirm the agreed HTTPS ingress and same-origin API behavior.}
\lead{Evidence log}{Save date, exact revision, build/config identifiers, environment, test fixture or manifest, pass/fail/blocked reason and owner. Record whether the live cause was proven or remains open. Keep secrets and private uploads out of evidence captures.}

\pagehead{Phase 7A}{Useful public site while reviewed findings are unavailable}
\label{firstpr}
\lead{First frontend PR}{Under the canonical app, edit \code{frontend/src/App.tsx}. Introduce a public layout in \code{components/layout/PublicLayout.tsx} with scoped styles. Under \code{pages/}, add \code{InvestigationPage.tsx} and \code{MethodsPage.tsx}. Update \code{RoadSafetyPage.tsx} and \code{services/roadSafetyService.ts} for explicit request states. Route the existing \code{HomePage.tsx} workspace entry through \code{/analyser}. Adjust workspace links that assumed root; preserve login return, protected deep links, existing dataset state and calculator URL. Keep this PR reviewable with no data-contract expansion.}
\lead{Root and Methods}{Use one county question, period, scope sentence, evidence status and named links. In the unavailable state, give a substantive “How the record was assembled” section with links to inclusion, coding and mapping. Methods explains WPRDC source and one-report unit, period and 2025 exclusion, selection rules, known-severity population and era boundary, missingness and coordinate screening, denominator/exposure limits, and citation guidance. Cite publisher, resource, period, extract/run, figure/table key and approved version where applicable. Static Methods loads independently of the data API; science editor verifies copy.}
\lead{Record deep link}{Keep \code{/road-safety} useful with the scope, Methods and status before any figure request. For a structured 503, say reviewed findings are unavailable. For network timeout, connection or unexpected transport failure, say the site could not check the findings. Offer retry, acknowledgement of progress and a route to Methods. Suppress numerical summaries and downloads while the gate is unavailable or cannot be established. Keep direct calculator and analyser access across the route split.}
\lead{Fixture and verification}{Reuse \code{backend/tests/overview_fixture.py}, \code{test_overview.py} and \code{frontend/tests/road-safety-smoke.mjs}. Exercise labeled synthetic, pending, missing, timeout/offline, recovery, and slow loading in controlled environments. Exclude fixtures from deployed images. Assert unique headings and landmarks, skip link, visible focus, menu state, no layout overflow at 390/768/1440 px and 200\% zoom, keyboard and screen-reader route/status discovery, direct refresh, back/forward, signed-in return and manual upload/export regression. Test public routes as anonymous. Capture root, Methods and record at wide and phone widths plus unavailable and transport states; record revision, browser, viewport, fixture, network state and traces.}
\lead{7A exit and rollback}{Frontend and science editor approve content, QA signs browser, reader and route regressions, deployment owner verifies intended HTTPS and affected access boundaries before exposure. Locally accepted PRs may wait for hosting. Keep a verified gate-closed image and route rollback. If deployment cannot safely expose the analyser, resolve its boundary before offering functional workspace access on that host.}

\pagehead{Phase 7B}{Populated investigation on verified HTTPS}
\lead{Objective and ownership}{Data steward and reviewer own exact-identity evidence and captions. Frontend and QA own one annual-record and selected-year coverage path. Deployment owner owns host, revalidation, rollback and sign-off. Use the 7A design and shared fixture states.}
\lead{First evidence slice}{At root lead with the approved annual plot, a bounded finding, caveat and source. Link to the annual record and year selector. Reuse Phase 6 charts and native table rendering where adequate. In \code{RoadSafetyPage.tsx} and the public figure component, join \code{county_periods} annual rows and \code{coordinate_coverage}/\code{field_coverage_by_year} using contract keys. Preserve the annual count, selected-year coverage and denominators in the chart, table, caption and data link. Identify gaps as approved observations and place the 2016 seam where relevant.}
\lead{API and row contract}{Read \code{/api/v1/road-safety/overview/provenance}, then verified \code{/tables/{table_id}}, \code{/tables/{table_id}/metrics/{column}}, \code{/downloads/{table_id}} and \code{/assets/{name}} as needed. Preserve kind, schema, version, manifest and per-output hash, run, units, geography, grain, ratio basis, input hashes, caveats and value state. Bind row annotations to complete keys in original response order before filtering. An annual key includes grain, period start and end. Use the same selection across chart/table/caption. Deep links carry the selected year and release context. Revalidate on open; explain withdrawn or replaced editions. Existing downloads remain checked original CSV bytes with provenance. Any scoped export includes key, population, filter and release identity, with approval checked at download time.}
\lead{Failure handling}{Abort superseded reads or ignore late responses by selection and release identity. A panel transport error with still-verified common provenance can retry locally. Gate failure, mixed identity or approval withdrawal clears protected views, share state, downloads and derived samples across mounted panels. Use no-store on protected data and assets. Revalidate before actions and on focus/return, with a bounded active-view interval agreed by deployment and reviewer. On ledger replacement verify the container sees new bytes. Rehearse withdrawal including ingress/cache behavior and stale tab responses.}
\lead{Exact release rehearsal}{Reviewer records a named, dated decision for this exact manifest in the real external ledger. Rehearse missing, pending, wrong-hash, tampered and changed-artifact cases with disposable fixtures before enabling any real release. From recorded held inputs, verify v6 checksum/lineage and reproduce D7 from a clean environment. Rerun applicable Phase 3/4/5 generation and refusal checks with documented skips, without promoting diagnostics. Reconcile annual, era, county and mapped candidate numbers before they appear publicly. Candidate anchors include 248,810 reports, 238,281 mapped, and 2024 coverage of 387 valid hours and zero populated illumination among 10,228 reports.}
\lead{Exit evidence}{Record labeled synthetic and genuine approved wide/phone captures, table and selection states, unavailable/mixed/retry/stale/withdrawal results, representative chart/table/export reconciliation, reader and accessibility results, redacted HTTPS and migration checks, backup/restore drill and signed runbook. Verify unsupported legacy site, EB, forecast, retrospective inference and real-treatment paths still refuse or remain unavailable. If approval or HTTPS evidence is missing, retain 7A and record 7B as pending. Roll back to a verified gate-closed deployment and clear protected state on discrepancy.}

\pagehead{Phase 7 design brief}{An Allegheny County civic research report}
\lead{Wide layout}{Use a masthead with Investigation, Methods, What-if calculator and Dataset analyser. Main uses a 60--72 character reading measure, with wider figures and a nearby caption/source margin where space allows. One display H1 (roughly 48--60 px), section H2s (28--36 px), readable 17--19 px prose, and discrete metadata. Use the title and a brief scope line, then the first approved figure when populated. When gated, put clear status and the Methods action before substantive explanatory sections.}
\lead{Phone layout}{Use a single column, 30--38 px title, 23--28 px H2, and readable body with generous line height. Keep visible “Menu” text or wrap text navigation, use accessible expanded state and focus order. Put status and Methods early, ahead of optional art; keep figure, caveat, source and table together in DOM order. Verify zoom and responsive data overflow without hiding data or headings.}
\lead{Tokens and type}{Candidate warm paper \texttt{\#F7F3E8}, ink \texttt{\#182021}, evidence teal \texttt{\#0E6059}, warning rust \texttt{\#8D3821}. Test actual contrast for body, muted/hover, focus, chart marks, visited links and status backgrounds. Use text plus shape/line styles for state and series. Candidate Instrument Serif → Georgia/serif, Source Sans 3 → system UI/Arial/sans-serif, IBM Plex Mono → ui-monospace/Consolas/monospace. Verify fallbacks and long hashes at phone width.}
\lead{Figure grammar}{Use a question or bounded takeaway, unit and grain, population, dates and geography, numerator and denominator for shares, exclusions/unknowns and nearby source/run link. Give a readable summary and a native table offering the same selection, keys, values and notes. Provide a keyboard selection control outside any pointer-dependent chart. Handle loading, unavailable and retry in text, announce dynamic changes without taking focus.}
\lead{Concept reference}{The supplied illustration is conceptual geography. Its generated copy is unreviewed. Use its ruled editorial hierarchy as a reference, move status ahead of optional art on mobile, and omit the illustration from the first populated view if it competes with the approved figure.}
\pagehead{Phase 7 design brief}{Public routes and workspace entry}
{\small
\begin{tabularx}{\linewidth}{@{}p{41mm}Y Y Y Y@{}}
\toprule Route & Reader question & First content and action & Evidence source & Unavailable state\\\midrule
\code{/} & What can this county record show? & Scope and first approved chart; read record or Methods. & Reviewed Methods and D7 annual rows. & Status, inclusion and measurement notes.\\
\code{/methods} & How were reports included and interpreted? & Scope and inclusion; jump to coding, coverage or citation. & Frozen v6 and reviewed source notes. & Methods and citations remain usable.\\
\code{/road-safety} & What was recorded by year? & Annual figure; select year or open table. & Approved annual and coverage rows. & Scope, distinct status, retry, Methods.\\
\code{/road-safety/map} & Where are reports concentrated? & Grid grain and scope; inspect the cells. & D7 grid and boundary. & Table if rendering fails; scope when gated.\\
\code{/road-safety/}\newline\code{evidence} & What do coding and coverage change? & Era and denominator; inspect comparable rows. & D7 missingness and associations. & Scope, status and Methods.\\
\code{/illustrative-}\newline\code{scenario} & What does a hypothetical factor change? & Example and explicit hypothetical form; calculate/reset. & Illustrative API and reference register. & Preserve inputs; actionable error and retry.\\
\code{/analyser} and \code{/login} & How do I analyse my own data? & Open workspace, learn from sample or sign in. & User data, approved aggregate sample or labeled teaching data. & Clear auth/scope state; safe teaching fallback.\\
\code{/road-safety/models} & What did retrospective evaluation show? & Feature timing and review scope; inspect approved comparisons. & Separately reviewed derivative or synthetic lesson. & Method note and eligible synthetic lesson.\\
\bottomrule
\end{tabularx}}
\sub{Layout boundaries}
Investigation and Methods share the reading shell, headings, figure and status styles. Keep the calculator focused on form, arithmetic and a result, with research navigation back to the report. Keep the analyser's workspace navigation, dense controls and project context under its own layout. Share type, focus treatment and selected identity tokens. Treat data provenance, hypothetical status and workspace ownership as separate labels.
\pagehead{Phase 7 · design review}{Reader validation and design review}
\lead{Evidence weighting}{Both UX reports support clear roles for public narrative, calculator and workspace, evidence near claims, recognizable controls and task-first testing. Controlled first-impression evidence addresses aesthetic ratings and familiarity. Comparisons of narrative and exploratory designs offer mixed results across recall, accuracy, preference and retrieval. Adjacent uncertainty and accessibility findings motivate caption and interaction choices, with direct EDAgent testing still needed. Practitioner case studies suggest patterns; community talk about interchangeable heroes, gradients and sparkles measures sentiment and cannot reliably identify authorship.}
\lead{Where the reports differ}{Gemini favors guided scrollytelling and suggests animated sensitivity; ChatGPT stresses reversible native reading order, qualified empirical findings and deterministic worked examples. Use native scroll and a static sensitivity range first. A user-specified range has no probability distribution. Treat Gemini's authentication and page-failure explanations as hypotheses where they exceed the audited source or screenshots. The reader cost of navigation framing, caveat placement, illustration, and future pacing remains open.}
{\small
\begin{tabularx}{\linewidth}{@{}p{32mm}Y Y@{}}
\toprule Marker & Better pattern & Assess on EDAgent\\\midrule
Generic copy & County, period and report grain in plain language. & Readers identify the subject and data unit unprompted.\\
Repetitive heroes/cards & Heading and section hierarchy, evidence where useful; familiar cards for genuine groups. & First-screen scan finds question, status and next action.\\
Decorative noise & Functional visual accents; subject-specific art only if useful. & Check figure salience, contrast, fallbacks and load with/without art.\\
Missing hierarchy & One primary action per reading step; contextual supporting links. & Trace title to figure, caveat and source on phone and wide.\\
Undisclosed limits & Denominator, count scope and coding/coverage near the claim. & Score caveat recall and zero-versus-missing interpretation.\\
Dead-end states & Distinct status, recovery and useful Methods. & Simulate network fail and pending release; complete next task.\\
Hover-only evidence & Labels, keyboard controls and native table. & Read selected row and caveat without a pointer.\\
Mobile order and focus & Early status, named menu, skip link, visible focus. & Check sequential keyboard/AT reading, zoom and reflow.\\
Weak chart alternative & Same units, coverage, data and selection in table. & Compare named chart values, table and eligible CSV.\\
\bottomrule
\end{tabularx}}
\lead{Formative study}{Aim for 12--16 county/public, civic-data and prospective analyser readers, including keyboard or assistive technology participants. Counterbalance variants and record device/order. Hold copy, data, links and styling constant while comparing editorial and panel layouts in the same state. Then hold layout fixed while varying substantive scope or caveat copy. Use plainly labeled fixtures until approval.}
\lead{Tasks and decisions}{Ask readers to find Methods and a source, state a finding and its limiting caveat, distinguish hypothetical results from observed effects, and find the analyser. Score accuracy, time, wrong turns, confidence and repeated misreadings. Collect recall and quality ratings before asking perceived authorship. Retest recurring or consequential errors. Weight comprehension, provenance and task completion above perceived authorship. Archive stimulus, rubric and decisions; palette, type, illustration and reading order remain hypotheses.}

\pagehead{Phase 8}{Map and severity evidence}
\lead{Objective and inputs}{Expose approved fixed-grain geographic and measurement evidence after the Phase 7 contract. Data steward supplies checked assets, science reviewer approves scope and chart language, frontend implements and QA validates chart/table parity. Fixture work can start early.}
\sub{8A · County grid and coverage}
In \code{frontend/src/pages/} add the map route and reuse the Figure and Overview service. Use D7 \code{historical_grid}, verified boundary and approved grid asset; check WGS 84, GEOID 42003, fixed full-period 0.001° cell keys and clipping. Candidate checks are 38,332 occupied cells summing to 238,281 mapped reports of 248,810. Reconcile usable coordinate, polygon exclusions and unmapped records. Show grid grain, coverage and historical concentration language ahead of controls. Keep precision, legend thresholds and mapped denominator consistent. Provide selectable cells and a keyboard-operable table or approved grid image with text and data fallback. Missing cells have no implied measured zero. Keep individual points, site ranking and route/address attribution out of this release.
\lead{Independent exit}{Verify asset hashes, totals, clipping and stable IDs, agreement of selected values, labeled fallback when map rendering fails, and gate closure when provenance fails. Capture map/table/failure at wide and phone widths and record phone load and interaction timing. Ask readers to explain the cell count and its coverage limit. Disable the map route or interactive layer if its verification fails.}
\sub{8B · Severity and measurement}
Build \code{/road-safety/evidence} from \code{severity_composition}, \code{source_missingness}, \code{severity_associations} and approved coverage rows. Render separate era panels, known and excluded populations, the 2016 seam and unsupported pooled nulls. Place the 2024 hour and illumination gaps beside the relevant fields. Introduce effect sizes with population, period and support (the recorded 2016--2023 association analysis uses 77,476 known-severity reports). Show sparse-level pooling, expected-cell checks and BH q-value state where applicable. Treat underflow as flagged null. Qualify as associations among reported crashes. New adjusted figures, joint distributions or imputed harmonisations need reviewed derivatives.
\lead{Independent exit}{Verify exact manifest, human-readable era distinction, denominator and exclusion reconciliation, chart/table/source agreement, keyboard operation and withdrawal. Score reader tasks on comparability and missingness. Save signed caption review and screenshots of available, partial and gated states. Expose map and evidence separately on verified HTTPS. Roll back each route independently to approved existing scope or Methods.}

\pagehead{Phase 9}{Illustrative calculator}
\lead{Objective and owners}{Frontend, backend and QA make Phase 5 arithmetic easy to run and hard to mistake for local evidence. Science editor reviews the treatment language. The contract and initial 7A reader feedback are sufficient to start; D7 approval and Phase 8 are independent.}
\lead{Work package}{In \code{InterventionPage.tsx} put “Run an invented example” before the manual form. Keep a focused tool layout with visible hypothetical status above inputs and in results. Use the same crash type and period for baseline and factor, retain decimal strings, schema 1 and server rounding. Show baseline × CMF, signed change and bounds with a common baseline. Example: 40 hypothetical crashes over 5 years × 0.80 = 32.00, reduction 8.00; bounds 0.70--0.95 yield 28.00--38.00. Label sensitivity as a user-set range without statistical confidence. Keep six \code{REFERENCE_ONLY} items separate from inputs and retain references and applicability caveats.}
\lead{State and safety}{Announce result and field errors; preserve valid inputs on recoverable transport failure, reset explicitly, and invalidate old output when inputs or request scope changes. Keep hypothetical status and assumptions with any export or share payload; reject or omit unsupported paths. Keep real-site refusal and factor compatibility checks in the backend interventions service.}
\lead{Independent exit}{Reuse \code{test_interventions.py} and browser smoke. Verify example arithmetic, neutral/increasing factor, zero baseline, malformed decimals, nonfinite/bounds, incompatible factor, stale/delayed requests and real-site refusal with no numbers. Capture initial, result, error and refusal states at desktop/phone. Reader completes an example, explains the sensitivity result and recovers from invalid input. Release on verified HTTPS independently, with fallback to the prior verified calculator.}

\pagehead{Phase 10}{General analyser}
\lead{Objective and owners}{Workspace/frontend maintains familiar task navigation and private data handling. Backend and privacy owner own isolation and lifecycle, QA owns account continuity. Require the repaired Phase 7 auth baseline. Retain existing account history and signed-in analysis/report behavior.}
\sub{10A · Onboarding and sample}
In \code{HomePage.tsx}, \code{UploadPage.tsx}, \code{WorkspacePage.tsx} and existing contexts, make the first manual task discoverable. Offer an allowlisted approved county aggregate sample with matching release identity, scope and withdrawal, or a clearly labeled invented teaching dataset. Keep the source, date and allowed interpretation with sample analyses and downloads. Verify sign-in return, sample-to-manual-analysis-to-export, workspace navigation and return to investigation. Exercise withdrawn/unavailable release and an independently usable teaching fallback. Capture fresh and signed-in paths and record task completion. Release onboarding/eligible sample after ownership, provenance and browser checks; turn off the county sample independently.
\sub{10B · Private guest lifecycle}
Use \code{AuthContext.tsx}, \code{DatasetContext.tsx}, \code{backend/app/core/deps.py}, existing dataset/session/job/report/export endpoints, repositories and storage. Use a distinct private guest principal; scope all records, files, jobs, downloads, caches and retries by owner. Set and communicate policy for expiry, upload/compute quotas, clear/delete, permitted formats, logging and account conversion before deployment. Use a migration only if needed. Treat repeated deletion as safe, stop or invalidate active work when revoked, prevent late results from recreating state and reclaim orphan files after restart. Provider transfer needs its own consent.
\lead{Independent exit}{Two isolated guests and one signed-in user each complete permitted upload, manual analysis and export. Regression verifies cross-owner denial including download links, missing login, expiry and deletion during work, restart, quota, failed upload, cleanup and existing account history. Review redacted lifecycle logs, browser captures and retention results at representative mobile/desktop widths, with keyboard and onboarding task checks. Rollback disables new guest creation while preserving account access and deletion/expiry obligations.}

\pagehead{Phase 11}{AI reliability, then optional BYOK}
\lead{Objective and owners}{AI/backend lead, privacy owner and QA separate connectivity, valid analytical plans, correct computation and grounded explanation. Phase 6 used mock mode. The teammate report is an unverified signal. Signed-in diagnostic work can start early. Guest support requires Phase 10 isolation.}
\sub{11A · Reliability and recovery}
Use \code{backend/app/api/v1/endpoints/chat.py}, existing config, orchestrator/LLM/context and grounding modules, \code{frontend/src/pages/ChatPage.tsx}, client and \code{test_chat_flow.py}. Verify observed runtime configuration, provider/model compatibility and a harmless fixed canary in the intended environment. Use synthetic tables with known operations, values, units and row counts. Assert approved operations and validated schemas, scoped session inputs and deterministic computation. Bind explanations to computed facts and IDs. Return the computed result and deterministic explanation when provider explanation fails. Keep manual analysis accessible and preserve signed-in account/session state.
\lead{Failure and release}{Exercise disabled, invalid credentials, rate limit, timeout, malformed/unsupported or overclaiming output, stale dataset switch and ambiguous retries. Define retry/latency budgets from measured traces, prevent duplicate computation or billing where feasible, and log stage/timing/error class without prompts, keys or data. Exit requires separate canary, full known-result and fallback evidence on tested environments, accessible status and recovery, and a short reader task from chat failure to manual answer. Routine rollback disables provider execution while preserving deterministic work.}
\sub{11B · Bring your own key (BYOK)}
Consider after 11A and a separate privacy and consent decision. Show provider, model, destination, data fields and lifecycle before transmission. Keep credentials request-scoped through proxy and provider call. Exclude them from persistent browser/server storage, logs, telemetry and exports. Provide stop/clear semantics and stale-response protection. Simulate invalid and revoked keys, reload, parallel users, session switch, cancellation and timeout using synthetic credentials and stub transport. Verify upstream rejection leaves application login and data intact. Run privacy and browser reviews and record eligibility scope. Release only with approved consent, credential isolation and lifecycle test evidence. Disable independently if any condition fails. Public D7 has no provider dependency.

\pagehead{Phase 12}{Retrospective model review}
\lead{Objective and owners}{ML/evaluation and data steward reproduce frozen work offline. Science reviewer controls exact-artifact publication; frontend and QA own reader understanding and gate behavior. Offline work can run alongside Phase 7.}
\sub{12A · Reproduction and methods}
Start from the frozen v6 notebook, \code{analysis/outputs/run_manifest.json}, \code{analysis/outputs/model_metrics/} and recorded scripts and environment. Reproduce with an independent runner using archived inputs. Record checksums, environments, comparator tolerances and discrepancies. Earlier audit records 15 CSVs within tolerance and five differing estimator digests. Preserve the original and qualify new results. Explain target (serious/fatal conditional on known recorded outcome), exclusions, feature timing and post-crash fields. Candidate post-2016 cohort is training 2016--2021 (58,660 known), test 2022--2024 (27,697 known, 1,231 positives); test years were previously inspected. Record split lineage and freeze rather than claim independent future validation. Exit with reproducible aggregate results or explicit discrepancies, scope and method review.
\sub{12B · Synthetic teaching}
If useful, add a deterministic dispersion lesson with synthetic mean/variance and seed if sampled. Use accessible data alternatives, persistent synthetic label, reset and labeled shared/exported view. Use no real-site inputs. Exit after deterministic reset, screen/keyboard checks and reader distinction between teaching and research evidence. Release independently of real model material.
\sub{12C · Exact-derivative review}
Propose offline aggregate builder and schema plus \code{/road-safety/models} only after acceptance. Key rows by cohort, split, feature set, model, calibration, threshold and approved grain. Include N, positives, holdout eligibility, timing, metric/uncertainty method and source/hash. Compare AP against prevalence and Brier against like-cohort references; keep ROC and AP distinct, document fit interval and calibration state, and identify collision ablation and post-event inputs. Show only stored, approved PR/ROC points, threshold/confusion counts or reliability bins with support. Do not infer curves from headline metrics or expose scores by case. Exit requires exact named review, reconciliation to frozen inputs, denominator/eligibility checks, browser and reader tests, and independent withdrawal test. Disable ungated comparison endpoints. Roll back to verified Methods and eligible synthetic lesson. Maintain all v3 public inference and site/treatment exclusions.

\pagehead{Release record}{Consolidated acceptance matrix}
{\small
\begin{tabularx}{\linewidth}{@{}p{22mm}Y Y@{}}
\toprule Gate & Pass evidence to archive & If open or failed\\\midrule
7 baseline & Target revision, host and config trace, strict identity and endpoint inventory, missing-ledger/refusal regressions. & Contain affected surfaces; continue isolated fixture work.\\
7A & Reviewed Methods, route checks, mobile/desktop captures, accessible states, task scores and intended HTTPS rehearsal. & Keep approved prior/safe gate-closed entry; retest defects.\\
7B & Named exact-manifest approval, independent reproduction, reconciled annual/coverage slice, target HTTPS and cache/revocation drill. & Clear protected data and serve 7A.\\
8A / 8B & Separate grid/coverage and era/association reconciliation, reader interpretation, verified fallback and withdrawal. & Disable affected route or artifact; retain scoped Methods.\\
9 & Hypothetical arithmetic, stale and refusal tests, reader task and responsive evidence. & Retain earlier verified calculator independently.\\
10A / 10B & Onboarding and sample identity checks, then isolation/expiry/deletion proof and signed-in regression. & Teaching fallback or disable guest entry; honor retention.\\
11A / 11B & Canary and known fixture with usable recovery; separate consent, isolation and request lifecycle. & Deterministic/manual mode, independent BYOK disable.\\
12A / B / C & Offline reproduction record, labeled synthetic reset, then exact derivative review and eligible cohort/caption checks. & Keep real output gated; show supported method or lesson.\\
\bottomrule
\end{tabularx}}
\sub{Evidence packet for each release}
Record revision and image, environment, checklist with dates and owners, test/fixture IDs, release hash and reviewer where applicable, screenshot and trace references, accessibility and reader issues, and rollback rehearsal. Keep scientific decisions and implementation sign-off separate. Log unresolved source gaps and the allowed exposure by route. For the Phase 7 handoff include a runnable demonstration of the approved slice and a separate gate-closed demonstration, hypothetical walkthrough, safe analyser path, limitations and an operational rollback guide.
\sub{V3 release obligations}
D0--D2 retain the fixed-input ledger, canonical entry, frozen execution and manifest. D3--D4 retain the aggregate scope and geographic conflict inventory with zero accepted site matches. D5 remains review-only. D6 stays closed. Phase 7 signs D7 exact-artifact release, D10 application, and D11 rehearsal; verify D8 source/status and D9 arithmetic in the existing simulator before sign-off. Phase 9 enhances it independently. Phase 12 extends D5 only through distinct review. Keep optional D12--D15 off the release path.
\sub{Source record}
\lead{Governing and handoff}{\lib{libfile_09c025eb71308191a73c321da0f26318}{V3 TeX} and \lib{libfile_919f39c5968c8191b2b6ac25646f7ae4}{v3 PDF}; \repo{blob/bcdf44d/docs/phase6_handoff.md}{Phase 6 handoff}; \repo{blob/93991e5/docs/phase3_phase4_review_status.md}{review status}; \repo{blob/bcdf44d/docs/descriptive_release_data_dictionary.md}{D7 API and dictionary}; \repo{commit/ca3197c}{directional pilot context}; \repo{commit/93991e5}{inspected main}. Review disposition comes from designated evidence, not commit titles.}
\lead{Issues and packages}{\lib{libfile_0cd101eb6b188191ae9378a5b07393c1}{EDAgent audit} and \lib{libfile_a6ec072a59c88191b85c6acb329bb373}{implementation plan A--H}, both at \code{87a7aea}. The audit is a revision-specific baseline. The commit and merge were accessible even though the named Phase 6 branch no longer resolved through GitHub.}
\lead{Design evidence}{\lib{libfile_069033a2f0a08191a055a614290e4e5c}{ChatGPT UX research} and \lib{libfile_fc9c1dc69af08191ad4fc186885ebd4f}{Gemini UX research}; \lib{libfile_115366c4a3b8819196cd824b1015c3f5}{Overview localhost}, \lib{libfile_7ee61624df50819192e52bd35efaa0e9}{calculator localhost} and \lib{libfile_bc3d7fe90d4c8191952197466a1992aa}{editorial concept}.}
\end{document}
